\documentclass[10pt]{article}

\usepackage[margin=0.75in]{geometry}
\usepackage[T1]{fontenc}
\usepackage{lmodern}
\usepackage{microtype}
\usepackage{amsmath,amssymb}
\usepackage[compact]{titlesec}
\usepackage{enumitem}
\usepackage[hidelinks]{hyperref}

\titleformat{\section}{\large\bfseries}{\thesection.}{0.5em}{}
\titlespacing*{\section}{0pt}{8pt}{3pt}
\titlespacing*{\paragraph}{0pt}{4pt}{0.6em}
\setlist{nosep,leftmargin=1.2em}
\setlength{\parskip}{3pt}
\setlength{\parindent}{0pt}
\setlength{\abovedisplayskip}{4pt}
\setlength{\belowdisplayskip}{4pt}
\newcommand{\code}[1]{\texttt{#1}}

\begin{document}

\noindent\textbf{Checkpoint 2: Adaptive Negotiation Training Through Human-Centered RL}
\hfill CS2822r, Fall 2026\\
{\small Yash Ravipati, Mahdi Nassereddine, Vincent Song}

% ---------------------------------------------------------------------------
\section{Research question}

Our core question is unchanged: \textbf{can rehearsing against an adaptive LLM
counterpart, followed by move-level feedback, improve how people negotiate in real
negotiations?} Building the MVP and the Checkpoint~1 feedback refined it in four ways.
\textbf{Scenario-specific rehearsal:} users upload the materials for the negotiation they
are about to have, and the system builds the case, a prep brief, and a training plan for
it. \textbf{Hybrid control:} the LLM writes the dialogue while code enforces the economics
(walk-away value and a concession schedule), because LLM counterparts concede too easily.
\textbf{Post-session counterfactual feedback:} we replay the turns that mattered with
alternative moves to estimate their value, with no live move suggestions, so the tool
builds skill rather than dependence. \textbf{A coach that plans over the user's state:}
the user is modeled as a hidden state of skills and errors with uncertainty, and a coach
uses it to choose what each session trains (Section~\ref{sec:change}).

\begin{itemize}
  \item \textbf{ML-Q1, steering:} do difficulty tiers move user score monotonically, and
    does a chosen focus create the situations it targets?
  \item \textbf{ML-Q2, user-state inference:} does the belief over a user's state match
    the truth for simulated users, with calibrated uncertainty?
  \item \textbf{ML-Q3, coaching:} does a coach with priors fitted across users fix new
    users' errors in fewer sessions than flat priors, our current curriculum, or no focus?
  \item \textbf{RQ1, transfer:} does rehearsal with counterfactual feedback raise real
    negotiation scores (z-scored within case and role) against a materials-only control?
  \item \textbf{RQ2, sim-to-real gap:} where does rehearsal behavior differ from real
    behavior, and how much preparation time does training need to transfer?
\end{itemize}

% ---------------------------------------------------------------------------
\section{What we have completed}

We have a working MVP covering all five workflow stages from Checkpoint~1: a local
Streamlit app (\code{negotrain}) with 93 automated tests. The flow is \emph{upload case
files $\rightarrow$ prep brief and training plan $\rightarrow$ rehearse $\rightarrow$
post-session feedback $\rightarrow$ profile and next scenario $\rightarrow$ log the real
outcome}, with everything stored locally.
\begin{itemize}
  \item \textbf{Case parser:} builds a structured case (issues, options, point tables,
    walk-away values) from PDF, Word, text, or ZIP uploads, and infers a labeled estimate
    if the counterpart's sheet is missing.
  \item \textbf{Prep brief and training plan:} an overview, the user's interests and
    walk-away, a checklist, and a warm-up $\rightarrow$ realistic $\rightarrow$ stretch
    sequence of rehearsals.
  \item \textbf{Hybrid counterpart:} Claude plays the other side in character, while code
    enforces its walk-away and concession schedule; five difficulty tiers set its
    anchoring, concession rate, cooperativeness, and aggression.
  \item \textbf{Rehearsal interface:} chat with formal offers, accept, walk away, and
    free-form custom terms (which are kept unscored rather than given invented points).
  \item \textbf{Counterfactual feedback:} after a session, the system finds the critical
    turns, replays each with alternative moves, and writes a coaching report on what each
    choice was worth.
  \item \textbf{Profiling, curriculum, and real-outcome log:} skill scores from behavioral
    features, a rule-based choice of the next rehearsal, and a log comparing real outcomes
    with rehearsals.
\end{itemize}
The counterpart never sees the user's confidential sheet, and the user sees the
counterpart's only in the debrief. The app runs through the Claude Code CLI, the
Anthropic API, or fully offline with a rule-based counterpart.

% ---------------------------------------------------------------------------
\section{How the tool has done}

\paragraph{In-class simulation.} We used an earlier version of the tool to prepare for the
in-class negotiation simulation. The main problem was time: the class gave us
15~minutes of preparation, which is not enough to train a negotiator. The tool is built
around repeated practice, in which the user rehearses, reads a debrief on what each move
was worth, and rehearses again against a counterpart adjusted to their weak spots. In
15~minutes a user can read a prep brief and perhaps start one rehearsal, but cannot
complete the rehearse--feedback--rehearse cycle where the training happens. The
simulation therefore told us more about how the tool fits into a real schedule than about
whether the training works: as built, it is a multi-session training program, not a
last-minute prep aid.

\paragraph{Algorithmic stress testing.} We ran 450 offline negotiations between the
rule-based counterpart (which shares the LLM counterpart's concession schedule) and
scripted users of three styles. The counterpart never agreed below its walk-away value,
and harder tiers reliably cost users points (Spearman $\rho = -0.82$ between tier and user
points). The tests also exposed three calibration problems: firm users never close at
tiers 2--5 because nothing acts as a deadline, the hardest tier pushes soft users all the
way down to their walk-away, and one difficulty setting (concession rate) does not affect
user points monotonically. \textbf{Limits:} apart from the class simulation, every
opponent so far has been scripted, and the profiler has no validation data yet.

% ---------------------------------------------------------------------------
\section{What we keep and what we change}
\label{sec:change}

\paragraph{Keep.} The hybrid counterpart (reliable economics, no leaks, difficulty that
works), scenario-specific rehearsal from uploaded files, post-session counterfactual
feedback, and strict information hiding.

\paragraph{Change 1: fit the tool to short prep windows.} The class simulation showed that
users may have only minutes before a negotiation. We will add a \emph{15-minute mode}: the
prep brief is generated the moment files are uploaded, followed by one short, capped
rehearsal aimed at the user's most costly likely error and a one-screen debrief. For
longer preparation, the tool will push users to start training days ahead, so that the
full rehearse--feedback cycle fits before the negotiation. RQ2 will test how much
preparation time training needs to transfer.

\paragraph{Change 2: a coach that plans over the user's state.} Following the Checkpoint~1
feedback, we replace the skill scores with a structured user state that carries
uncertainty. After session $t$ the user is a hidden vector $z_t = (k_t, \rho_t, \theta_t)$:
awareness of the case's key facts, rapport, and error rates such as weak anchoring,
conceding first, or accepting a low deal. Code counts each error $j$ per opportunity
($x_{j,t}$ errors in $o_{j,t}$ chances) and keeps a Beta belief with forgetting, so the
belief notices improvement; one Claude call per session tags facts and rapport. Before
each session the coach samples the belief and picks the focus $f_t$ with the largest
expected gain (Thompson sampling), where $w_j$ is the points error $j$ costs (from the
counterfactual replays) and $\delta_j$ is how much one focused session reduces it:
\begin{equation}
\theta_j \sim \mathrm{Beta}\left(\alpha_j + X_j,\; \beta_j + O_j - X_j\right),\quad
X_j \leftarrow \lambda X_j + x_{j,t},\quad O_j \leftarrow \lambda O_j + o_{j,t};
\qquad
f_t = \arg\max_j\, w_j\, \tilde{\theta}_j\, \delta_j .
\end{equation}
The focus steers the counterpart (for example, one that withholds its first concession to
test conceding first), the debrief, and the prep checklist. Because each user has only a
few sessions, the priors $\phi = (\alpha_j, \beta_j, \delta_j, \lambda)$ are fitted across
users, first by empirical Bayes and then by direct optimization of how much users improve,
in the setting of \href{https://proceedings.mlr.press/v139/kveton21a.html}{Meta-Thompson
Sampling} (Kveton et al., 2021). Code makes the decisions and Claude writes the words, so
we can train on thousands of free simulated users and validate on 10--20 Claude-played
users (about \$150--300) before the human pilot.

\paragraph{Change 3: calibration and evidence.} Add a visible deadline so firm negotiators
close, keep the hardest tier for users who have earned it, and fix the concession-rate
setting. Use 4--5 rollouts per alternative with confidence intervals in the debrief. Add a
human-vs-human mode so real negotiations can be run and logged in the tool.

\paragraph{Next steps.} We build the 15-minute mode, the user-state belief, and the coach
first, since they directly improve the tool; then run the steering, inference, and
coaching experiments in simulation; and in parallel prepare a mid-November pilot with
12--20 participants, a treatment group that trains with the tool against a materials-only
control, and five real negotiations each.

\end{document}
